\specialchapter{根系的主要性质概述}

（在本概述中，我们限于约化根系的情形，并在实数域上工作。）

\begin{enumerate}
\def\labelenumi{\arabic{enumi})}
\tightlist
\item
  令 V 为实向量空间。V 中的约化根系是 V 的一个子集 R，具有以下性质：
\end{enumerate}

\begin{enumerate}
\def\labelenumi{(\roman{enumi})}
\item
  R 是有限集并且生成 V。
\item
  对所有 \(\alpha \in \mathbb{R}\)，存在 \(\alpha^{\vee} \in V^{*}\)，使得 \(\langle \alpha, \alpha^{\vee} \rangle = 2\)，并且使映射
\end{enumerate}

\[
s _ {\alpha}: x \mapsto x - \langle x, \alpha^{\vee} \rangle \alpha
\]

从 V 到 V 将 R 映到 R。

\begin{enumerate}
\def\labelenumi{(\roman{enumi})}
\setcounter{enumi}{2}
\item
  对所有 \(\alpha \in \mathbb{R}\)，有 \(\alpha^{\vee}(\mathbb{R}) \subseteq \mathbf{Z}\)。
\item
  若 \(\alpha \in \mathbb{R}\)，则 \(2\alpha \notin \mathbb{R}\)。
\end{enumerate}

鉴于 (i)，由 (ii) 保证存在的元素 \(\alpha^{\vee}\) 是唯一的；因此 (iii) 是有意义的。映射 \(s_{\alpha}\) 是一个反射，它固定 \(\mathrm{L}_{\alpha} = \mathrm{Ker}(\alpha^{\vee})\) 中的点，并将 \(\alpha\) 变为 \(-\alpha\)。

R 的元素称为根。V 的维数称为根系的秩。

\begin{enumerate}
\def\labelenumi{\arabic{enumi})}
\setcounter{enumi}{1}
\item
  保持 R 不变的 V 的自同构群记为 A(R)。\(s_{\alpha}\)（\(\alpha \in R\)）生成 A(R) 的一个子群 W(R)，称为 R 的 Weyl 群；该子群在 A(R) 中正规。属于 W(R) 的唯一反射是 \(s_{\alpha}\)。
\item
  集合 \(R^{-}\) 由 \(\alpha^{-}\)（其中 \(\alpha \in R\)）组成，是 \(V^{*}\) 中的一个约化根系，称为 R 的对偶根系。映射 \(\alpha \mapsto \alpha^{-}\) 是从 R 到 \(R^{-}\) 的一个双射，称为典范双射。我们有 \((\mathrm{R}^{-})^{-} = \mathrm{R}\)，并且典范双射 \(R \to R^{-}\) 和 \(R^{-} \to R\) 互为逆映射。映射 \(u \mapsto {}^{t}u^{-1}\) 给出从 \(W(R)\) 到 \(W(R^{-})\) 的一个同构；借助该同构，我们认同这两个群。
\item
  令 V 为实向量空间，并且是向量子空间 \(V_{1},\ldots,V_{r}\) 的直和。对每个 i，令 \(R_{i}\) 为 \(V_{i}\) 中的约化根系。则各 \(R_{i}\) 的并集 R 是 V 中的一个根系，称为各 \(R_{i}\) 的直和。群 \(W(R)\) 可认同为各 \(W(R_{i})\) 的直积。若 \(R \neq \varnothing\)，且 R 不是两个非空根系的直和，则称 R 是不可约的。这等价于说 \(W(R)\) 是不可约的。每个约化根系
\end{enumerate}

R 是不可约约化根系的直和；这些分支除排列外唯一确定，并称为 R 的不可约分支。

\begin{enumerate}
\def\labelenumi{\arabic{enumi})}
\setcounter{enumi}{4}
\item
  令 R 为 V 中的约化根系。V 上存在在 \(W(R)\) 作用下不变的标量积。下文用 \((x|y)\) 表示这样的标量积。若将 V 与 \(V^{*}\) 用 \((x|y)\) 认同，则有 \(\alpha^{\vee} = \frac{2\alpha}{(\alpha|\alpha)}\)。反射 \(s_{\alpha}\) 是将 \(\alpha\) 变为 \(-\alpha\) 的正交反射。若 R 不可约，则 Weyl 群在给定长度的根的集合上传递作用。若 R 不可约，则标量积 \((x|y)\) 除相差一个常数因子外是唯一的。
\item
  令 R 为一个约化根系。对 \(\alpha, \beta \in \mathbb{R}\)，置
\end{enumerate}

\[
\langle \alpha , \beta^{\vee} \rangle = n (\alpha , \beta) \in \mathbf {Z}.
\]

我们有

\[
\begin{array}{c} n (\alpha , \alpha) = 2, \\ s _ {\beta} (\alpha) = \alpha - n (\alpha , \beta) \beta , \\ n (\alpha , \beta) = \frac {2 (\alpha | \beta)}{(\beta | \beta)}. \end{array}
\]

除交换 \(\alpha\) 与 \(\beta\) 外，只有以下可能性：

\[
n (\alpha , \beta) = n (\beta , \alpha) = 0; \quad (\alpha , \beta) = \frac {\pi}{2}; \quad s _ {\alpha} s _ {\beta} \text {   of   order   } 2;
\]

\[
n (\alpha , \beta) = n (\beta , \alpha) = 1; \quad (\alpha , \beta) = \frac {\pi}{3}; \quad s _ {\alpha} s _ {\beta} \text {   of   order   } 3;
\]

\[
n (\alpha , \beta) = n (\beta , \alpha) = - 1; \quad (\alpha , \beta) = \frac {2 \pi}{3}; \quad s _ {\alpha} s _ {\beta} \text {   of   order   } 3;
\]

\[
n (\alpha , \beta) = 1, n (\beta , \alpha) = 2; \quad (\alpha , \beta) = \frac {\pi}{4}; \quad s _ {\alpha} s _ {\beta} \text {   of   order   } 4;
\]

\[
n (\alpha , \beta) = - 1, n (\beta , \alpha) = - 2; (\alpha , \beta) = \frac {3 \pi}{4}; s _ {\alpha} s _ {\beta} \text {   of   order   } 4;
\]

\[
n (\alpha , \beta) = 1, n (\beta , \alpha) = 3; \quad (\alpha , \beta) = \frac {\pi}{6}; \quad s _ {\alpha} s _ {\beta} \text {   of   order   } 6;
\]

\[
n (\alpha , \beta) = - 1, n (\beta , \alpha) = - 3; (\widehat {\alpha , \beta}) = \frac {5 \pi}{6}; s _ {\alpha} s _ {\beta} \text {   of   order   } 6;
\]

\[
n (\alpha , \beta) = n (\beta , \alpha) = 2; \quad \alpha = \beta ; \quad n (\alpha , \beta) = n (\beta , \alpha) = - 2; \quad \alpha = - \beta ;
\]
\[
\begin{gathered}
 n (\alpha , \beta) = n (\beta , \alpha) = 1 \Longrightarrow \| \alpha \| = \| \beta \| ; \quad n (\alpha , \beta) = n (\beta , \alpha) = -1 \Longrightarrow \| \alpha \| = \| \beta \| ; \\
 (n (\alpha , \beta) , n (\beta , \alpha)) = (1,2) \Longrightarrow \| \beta \| = \sqrt {2} \| \alpha \| ; \quad (n (\alpha , \beta) , n (\beta , \alpha)) = (-1,-2) \Longrightarrow \| \beta \| = \sqrt {2} \| \alpha \| ; \\
 (n (\alpha , \beta) , n (\beta , \alpha)) = (1,3) \Longrightarrow \| \beta \| = \sqrt {3} \| \alpha \| ; \quad (n (\alpha , \beta) , n (\beta , \alpha)) = (-1,-3) \Longrightarrow \| \beta \| = \sqrt {3} \| \alpha \| .
\end{gathered}
\]

\begin{enumerate}
\def\labelenumi{\arabic{enumi})}
\setcounter{enumi}{6}
\item
  令 \(\alpha, \beta \in \mathbb{R}\)。若 \((\alpha|\beta) > 0\)，则 \(\alpha + \beta\) 是一个根，除非 \(\alpha = \beta\)。若 \((\alpha|\beta) < 0\)，则 \(\alpha + \beta\) 是一个根，除非 \(\alpha = -\beta\)。
\item
  令 \(\alpha, \beta\) 为两个不成比例的根。集合 I 由 \(j \in \mathbf{Z}\) 组成；其中 \(\beta + j\alpha \in \mathbb{R}\) 成立。它是包含 0 的、记为 \([-q, p]\) 的 \(\mathbf{Z}\) 区间。我们有
\end{enumerate}

\[
p - q = - n (\beta , \alpha), \quad \frac {q + 1}{p} = \frac {(\beta + \alpha | \beta + \alpha)}{(\beta | \beta)}.
\]

令 S 为 \(\beta + j\alpha\)（\(j \in I\)）所成的集合。则 \(s_{\alpha}(S) = S\) 且 \(s_{\alpha}(\beta + p\alpha) = \beta - q\alpha\)。我们称 S 为 \(\alpha\)-根链，它由 \(\beta\) 定义；称 \(\beta - q\alpha\) 为其起点，称 \(\beta + p\alpha\) 为其终点，并称 \(p + q\) 为其长度。

若 T 是 \(\alpha\)-根链，起点为 \(\gamma\)，则 T 的长度为 \(-n(\gamma, \alpha)\)。

\begin{enumerate}
\def\labelenumi{\arabic{enumi})}
\setcounter{enumi}{8}
\item
  令 X 为 \(\operatorname{Ker}\alpha^{\vee}\)（\(\alpha\in R\)）的并集。V-X 的连通分支称为 R 在 V 中的室。它们是开单纯锥。Weyl 群在室的集合上单传递地作用。若 C 是一个室，则 \(\overline{C}\) 是 \(\mathrm{W}(\mathrm{R})\) 的基本域。我们有 \((x|y)>0\) 对 \(x,y\in C\) 成立。由从 V 到 \(V^{*}\)、与 \((x|y)\) 对应的双射，将 V 中 R 的室集合双射到 \(R^{\vee}\) 中 \(V^{*}\) 的室集合；我们把 \(C^{\vee}\) 记为 C 在此双射下的像。
\item
  令 C 为 R 的一个室。令 \(L_{1}, L_{2}, \ldots, L_{l}\) 为 C 的壁。对每个 i，存在唯一的根 \(\alpha_{i}\)，使得 \(L_{i} = L_{\alpha_{i}}\)，且 \(\alpha_{i}\) 与 \(L_{i}\) 位于同一侧。族 \((\alpha_{1}, \ldots, \alpha_{l})\) 是 V 的一个基，而 C 是 \(x \in V\) 中满足对所有 i 有 \(\langle \alpha_{i}, x \rangle > 0\) 的元素所成的集合，换言之，是满足对所有 i 有 \((\alpha_{i} | x) > 0\) 的集合。于是 \(\{\alpha_{1}, \ldots, \alpha_{l}\}\) 称为由 C 定义的 R 的基 B(C)。我们有 \((\alpha_{i} | \alpha_{j}) \leqslant 0\)；当 \(i \neq j\) 时。Weyl 群在基的集合上单传递地作用。每个根都可由 W(R) 的某个元素变换为 B(C) 中的一个元素。我们有 \(\{\alpha_{1}^{\vee}, \ldots, \alpha_{l}^{\vee}\} = B(C^{\vee})\)。
\item
  置 \(s_{\alpha_i} = s_i\)，令 \(S\) 为各 \(s_i\) 所成的集合，并令 \(m_{ij}\) 为 \(s_i s_j\) 的阶。对 (W(R), S) 而言，这是带有 Coxeter 矩阵（\(m_{ij}\)）的 Coxeter 系统；换言之，W(R) 由生成元（\(s_i\)）\(_{1 \leqslant i \leqslant l}\) 及关系（\(s_i s_j\)）\(^{m_{ij}} = 1\) 定义。W(R) 中的两个元素 \(s_i\) 和 \(s_j\) 共轭，当且仅当存在指标序列（\(i_1, i_2, \ldots, i_q\)），使得 \(i_1 = i, i_q = j\)，并且每个 \(m_{i_l i_{l+1}}\) 都等于 3。
\item
  令 \(n_{ij} = n(\alpha_i, \alpha_j)\)。矩阵 \((n_{ij})_{1 \leqslant i,j \leqslant l}\) 称为 \(R\) 的 Cartan 矩阵。它（除 \(1, 2, \ldots, l\) 的排列外）与 \(C\) 的选择无关。我们有 \(n_{ii} = 2\)。\(n_{ij} \in \{0, -1, -2, -3\}\) 对 \(i \neq j\) 成立。如果两个根系具有相同的 Cartan 矩阵，则它们同构。
\item
  令 \(G\) 为 \(A(R)\) 中保持 \(B(C)\) 不变的子群。则 \(A(R)\) 是 \(G\) 与 \(W(R)\) 的半直积。
\item
  由 C 在 V（分别为 \(V^{*}\)）上定义的序关系，是与 V（分别为 \(V^{*}\)）的向量空间结构相容的序关系；其中 \(\geqslant 0\) 的元素是 \(\alpha_{i}\)（分别为 \(\alpha_{i}^{\vee}\)）的系数为 \(\geqslant 0\) 的线性组合。这些元素称为关于 C 或 B(C) 的正元素。这些序关系也由 \(C^{\vee}\) 定义。V 的一个元素为 \(\geqslant 0\)，当且仅当它在 \(C^{\vee}\) 上的取值为 \(\geqslant 0\)。关于 C 的 \(\geqslant 0\) 元素集合包含 \(\overline{C}\)，但一般不同于 \(\overline{C}\)。令 \(x \in V\)。则 \(x \in \overline{C}\) 当且仅当 \(x \geqslant w(x)\) 对所有 \(w \in W(R)\) 都成立；而 \(x \in C\) 当且仅当 \(x > w(x)\) 对所有不同于 1 的 \(w \in W(R)\) 都成立。
\item
  每个根对于 \(C\) 不是正的就是负的。我们把 \(R_{+}(C)\) 记为对于 \(C\) 为正的根所成的集合，于是 \(R = R_{+}(C) \cup (-R_{+}(C))\) 是 R 的一个划分。反射 \(s_{i}\) 将 \(\alpha_{i}\) 变为 \(-\alpha_{i}\)，并将 \(\mathrm{R}_{+}(\mathrm{C})\) 中异于 \(\alpha_{i}\) 的元素彼此置换。
\item
  设 B 是 R 的一个基。对于 B 为正的（分别为负的）每个根，都是 B 中元素的线性组合，其系数全为整数 \(\geqslant 0\)（分别为 \(\leqslant 0\)）。
\item
  设 \((\beta_{1},\beta_{2},\ldots,\beta_{n})\) 是一列对于 C 为正的根，并且 \(\beta_{1}+\cdots+\beta_{n}\) 是一个根。存在一个置换 \(\pi\in S_{n}\)，使得对所有 \(i\in\{1,2,\ldots,n\}\)，\(\beta_{\pi(1)}+\beta_{\pi(2)}+\cdots+\beta_{\pi(i)}\) 都是一个根。
\item
  设 \(\alpha \in \mathrm{R}_{+}(\mathrm{C})\)。则 \(\alpha \in \mathrm{B}(\mathrm{C})\)，当且仅当 \(\alpha\) 不能写成正根之和。
\item
  设 C 为一个室，\((\alpha_{1},\alpha_{2},\ldots,\alpha_{l})\) 为相应的基。对 \(I=\{1,2,\ldots,l\}\) 的任意子集 J，令 \(W_{J}\) 为 \(W(R)\) 的由 \(s_{i}\)（满足 \(i\in J\)）生成的子群。令 \(C_{J}\) 为 \(\alpha_{j}\)（\(j\in J\)）的系数为 \textgreater0 的线性组合所成的集合，于是 \(C_{J}\) 是 C 的一个面。
\end{enumerate}

令 \(J \subset I\)，\(g \in W(R)\)。以下条件等价：

\begin{enumerate}
\def\labelenumi{\alph{enumi})}
\item
  g 使 \(C_{J}\) 中的某一点不变：
\item
  g 使 \(C_{J}\) 中的每一点不变；
\item
  \(g\) 使 \(\overline{\mathbf{C}_{\mathrm{J}}}\) 中的每一点不变；
\end{enumerate}

\(d) g(C_{J}) = C_{J};\)

\begin{enumerate}
\def\labelenumi{\alph{enumi})}
\setcounter{enumi}{4}
\tightlist
\item
  \(g(\overline{\mathrm{C}_{\mathrm{J}}}) = \overline{\mathrm{C}_{\mathrm{J}}}\) ;
\end{enumerate}

\(f) g \in W_{J}\) .

令 \(J, J' \subset \{1, 2, \ldots, l\}\)，且 \(g, g' \in W(R)\)。以下条件等价：

\begin{enumerate}
\def\labelenumi{\alph{enumi})}
\item
  \(g(\mathrm{C}_{\mathrm{J}}) = g'(\mathrm{C}_{\mathrm{J}'})\) ;
\item
  \(g(\mathrm{C}_{\mathrm{J}})\cap g'(\mathrm{C}_{\mathrm{J}'})\neq \emptyset ;\)
\item
  \(g\mathrm{W}_{\mathrm{J}} = g^{\prime}\mathrm{W}_{\mathrm{J}^{\prime}}\)
\end{enumerate}

\(d) J = J' \text{ 且 } g' \in gW_J.\)

令 \(J_1, J_2, \ldots, J_r \subset I\)，且 \(J = J_1 \cap \cdots \cap J_r\)。则 \(W_J = W_{J_1} \cap \cdots \cap W_{J_r}\)。对所有 \(g \in W(R)\)，存在 \(J \subset I\)，使得 \(\overline{C} \cap g(\overline{C}) = \overline{C}_J\) 且 \(g \in W_J\)。

\begin{enumerate}
\def\labelenumi{\arabic{enumi})}
\setcounter{enumi}{19}
\tightlist
\item
  令 \(\mathbf{P}\) 为 \(\mathbb{R}\) 的一个子集。称 \(\mathbf{P}\) 为闭的，是指条件 \(\alpha \in \mathbb{P}, \beta \in \mathbb{P}, \alpha + \beta \in \mathbb{R}\) 蕴含 \(\alpha + \beta \in \mathbb{P}\)；称 \(\mathbf{P}\) 为抛物的，是指 \(\mathbf{P}\) 闭且满足 \(\mathbf{P} \cup (-\mathbf{P}) = \mathbb{R}\)。以下条件等价：
\end{enumerate}

\begin{enumerate}
\def\labelenumi{\alph{enumi})}
\item
  P 是抛物集；
\item
  P 是闭集，并且存在一个室 C，使得 \(P \supset R_{+}(C)\)；
\item
  存在一个室 C 和 B(C) 的一个子集 \(\Sigma\)，使得 P 是 \(R_{+}(C)\) 与根的集合 Q 的并集，其中 Q 由 \(\Sigma\) 中元素的系数为整数 \(\leqslant 0\) 的线性组合组成。
\end{enumerate}

假设这些条件成立，并令 \(V_{1}\) 为 V 中由 \(\Sigma\) 生成的向量子空间。则

\[
P \cap (- P) = Q \cup (- Q) = V _ {1} \cap R,
\]

且 \(P \cap (-P)\) 是 \(V_1\) 中以 \(\Sigma\) 为基的根系。

令 \(P', C', \Sigma'\) 具有相应的性质。若存在 \(W(R)\) 的一个元素将 P 变为 \(P'\)，则存在 \(W(R)\) 的一个元素，将 C 变为 \(C'\)，将 \(\Sigma\) 变为 \(\Sigma'\)，并将 P 变为 \(P'\)。

\begin{enumerate}
\def\labelenumi{\arabic{enumi})}
\setcounter{enumi}{20}
\tightlist
\item
  令 \(\mathbf{P}\) 为 \(\mathbb{R}\) 的一个子集。以下条件等价：
\end{enumerate}

\begin{enumerate}
\def\labelenumi{\alph{enumi})}
\item
  存在一个室 C，使得 \(P = R_{+}(C)\)：
\item
  P 是闭集，且 \(\{P, -P\}\) 是 R 的一个划分。
\end{enumerate}

此时室 C 是唯一的。

假设 V 配备了有序向量空间的结构，使得每个根或者为正或者为负。令 \(R_{+}\) 为关于此结构的正根所成的集合。存在唯一的室 C，使得 \(\mathrm{R}_{+} = \mathrm{R}_{+}(\mathrm{C})\)。

\begin{enumerate}
\def\labelenumi{\arabic{enumi})}
\setcounter{enumi}{21}
\item
  R 的子集 B 是 R 的一个基，当且仅当 B 中的元素线性无关，并且每个根都是 B 中元素的线性组合，且其系数要么全为 \(\geqslant 0\)，要么全为 \(\leqslant 0\)。
\item
  令 \(\mathbf{P}\) 为 \(\mathbf{R}\) 的一个闭子集，并满足 \(\mathbf{P} \cap (-\mathbf{P}) = \varnothing\)。存在一个室 \(\mathbf{C}\)，使得 \(\mathbf{P} \subset R_{+}(\mathbf{C})\)。
\item
  若 P = -P，则称 R 的子集 P 为对称的。令 P 为 R 的一个子集，并令 \(V_{1}\)（分别为 \(\Gamma\)）为 V 中由 P 生成的向量子空间（分别为加法子群）。以下条件等价：
\end{enumerate}

\begin{enumerate}
\def\labelenumi{\alph{enumi})}
\item
  P 是闭的且对称的；
\item
  P 是闭的，并且是 \(V_{1}\) 中的一个根系；
\item
  \(\Gamma \cap \mathbb{R} = \mathbb{P}\) .
\end{enumerate}

\begin{enumerate}
\def\labelenumi{\arabic{enumi})}
\setcounter{enumi}{24}
\tightlist
\item
  假设 R 不可约。令 C 为一个室，置
\end{enumerate}

\[
\mathrm{B} (\mathrm{C}) = \left\{\alpha_ {1}, \dots , \alpha_ {l} \right\}.
\]

R 中存在最高元素（相对于 C 所定义的序），即存在 R 的一个元素 \(\tilde{\alpha}=n_{1}\alpha_{1}+\cdots+n_{l}\alpha_{l}\)，使得对于任意根

\[
p _ {1} \alpha_{1} + \dots + p _ {l} \alpha_{l}.
\]

\(n_1 \geqslant p_1, \ldots, n_l \geqslant p_l\)。于是 \(\tilde{\alpha} \in \overline{C}\)，并且有 \(\| \tilde{\alpha} \| \geqslant \| \alpha \|\)（对每个根 \(\alpha\)）。

\begin{enumerate}
\def\labelenumi{\arabic{enumi})}
\setcounter{enumi}{25}
\tightlist
\item
  V 中由 R 生成的子群记为 Q(R)；Q(R) 的元素称为 R 的根权。群 Q(R) 是 V 中秩为 \(l = \dim V\) 的离散子群。R 的任一基都是 Q(R) 的一个基。
\end{enumerate}

与 \(Q(R^{\vee})\) 对应的 V 的子群记为 \(P(R)\)；\(P(R)\) 的元素称为 R 的权。群 \(P(R)\) 是 V 中秩为 l、包含 \(Q(R)\) 的离散子群。群 \(P(R)/Q(R)\)、\(P(R^{\vee})/Q(R^{\vee})\) 是有限群且彼此同构；它们的阶 f 称为 R 的连接指标。按 25) 的记号，\(W(R)\) 的阶为 \(l!n_{1}n_{2}\ldots n_{l}f\)。

群 \(A(R)\) 保持 \(P(R)\)、\(Q(R)\) 不变，因而作用于 \(P(R)/Q(R)\)。群 \(W(R)\) 在 \(P(R)/Q(R)\) 上作用平凡，所以 \(A(R)/W(R)\) 作用于 \(P(R)/Q(R)\)。

\begin{enumerate}
\def\labelenumi{\arabic{enumi})}
\setcounter{enumi}{26}
\tightlist
\item
  令 C 为一个室。令 \(B = (\alpha_{1}, \ldots, \alpha_{l})\) 为相应的 R 的基。对偶基 \((\omega_{1}, \ldots, \omega_{l})\)（相对于 \((\alpha_{1}^{\vee}, \ldots, \alpha_{l}^{\vee})\)）是 P(R) 的一个基。称 \(\omega_{i}\) 为基本权（关于 C 或关于 B）。\(\omega_{i}\) 的系数为 \textgreater{} 0（分别为 \(\geqslant 0\)）的线性组合所成的集合是 C（分别是 \(\overline{C}\)）。\(\omega_{i}\) 的系数为整数 \(\geqslant 0\) 的线性组合称为支配权。P(R) 的每个元素都可由 W(R) 的某个元素变换为唯一的支配权。支配权正是 V 中的元素 \(\omega\)，它满足 \(\frac{2(\omega|\alpha_{i})}{(\alpha_{i}|\alpha_{i})}\) 是一个 \(\geqslant 0\) 的整数。
\end{enumerate}

\[
28) \text {   Let   } \rho = \frac {1}{2} \sum_ {\alpha \in R _ {+} (C)} \alpha . \text {   Then,   } \rho = \omega_ {1} + \dots + \omega_ {l} \in C.
\]

\begin{enumerate}
\def\labelenumi{\arabic{enumi})}
\setcounter{enumi}{28}
\item
  令 \(\mathbf{T}\) 为 \(\mathbf{V}^*\) 的平移群，其向量属于 \(\mathrm{Q}(\mathbb{R}^{\vee})\)。\(\mathbf{V}^*\) 的仿射变换群由 \(\mathbf{T}\) 和 \(\mathrm{W}(\mathbb{R})\) 生成，是 \(\mathrm{W}(\mathbb{R})\) 与 \(\mathbf{T}\) 的半直积。该群称为 \(\mathbb{R}\) 的仿射 Weyl 群，记为 \(\mathrm{W}_a(\mathbb{R})\)。它在 \(\mathbf{V}^*\) 上适当作用。对 \(\alpha \in \mathbb{R}\) 和 \(\lambda \in \mathbf{Z}\)，令 \(s_{\alpha, \lambda}\) 为映射 \(x^* \mapsto x^* - \langle x^*, \alpha \rangle \alpha^{\vee} + \lambda \alpha^{\vee}\)；这是一个仿射反射，而其不动点集 \(L_{\alpha, \lambda}\) 由方程 \(\langle x^*, \alpha \rangle = \lambda\) 定义；我们有 \(L_{\alpha, \lambda} = L_{\alpha} + \frac{1}{2}\lambda \alpha^{\vee}\)。\(s_{\alpha, \lambda}\) 是属于 \(W_a(\mathbb{R})\) 的仿射反射，并且它们生成群 \(W_a(\mathbb{R})\)。
\item
  令 \(\mathbf{E}\) 为各 \(L_{\alpha, \lambda}\) 的并集，其中 \(\alpha \in \mathbb{R}\) 且 \(\lambda \in \mathbf{Z}\)。\(V^{*} - E\) 的连通分支称为 \(\mathbb{R}\) 的小室。若 \(\mathbb{R}\) 不可约，则每个小室都是开单纯形；一般地，小室是开单纯形的乘积。群 \(W_{a}(\mathbb{R})\) 在小室的集合上单传递地作用。若 \(\mathbb{C}\) 是一个小室，则 \(\overline{\mathbb{C}}\) 是 \(W_{a}(\mathbb{R})\) 的基本域。令 \(\sigma_{1}, \ldots, \sigma_{q}\) 为 \(W_{a}(\mathbb{R})\) 中与 \(\mathbb{C}\) 的壁对应的反射；令 \(\mu_{ij}\) 为 \(\sigma_{i}\sigma_{j}\) 的阶。则 \(W_{a}(\mathbb{R})\) 由生成元 \(\sigma_{i}\) 及关系 \((\sigma_{i}\sigma_{j})^{\mu_{ij}} = 1\) 定义。
\item
  若 \(p \in \mathrm{P}(\mathbb{R}^{\vee})\)，则存在一个小室 \(C\)，使得 \(p\) 是 \(\overline{\mathbb{C}}\) 的一个极点。若 \(\mathbf{C}'\) 是一个小室，则存在唯一的根权，它是 \(\overline{\mathbb{C}'}\) 的一个极点。
\end{enumerate}

令 \(x^{*} \in V^{*}\)；以下条件等价：

\begin{enumerate}
\def\labelenumi{\alph{enumi})}
\item
  \(x^{*} \in \mathrm{P}(\mathrm{R}^{\vee})\) ;
\item
  对所有 \(\alpha \in \mathbb{R}\)，与 \(\mathbf{L}_{\alpha}\) 平行且经过 \(x^{*}\) 的超平面是 \(\mathbf{L}_{\alpha,\lambda}\) 中的一个。
\end{enumerate}

令 \(C'\) 为 \(R^{\vee}\) 的一个室。存在唯一的小室 \(C\) 包含于 \(C'\) 中，并满足 \(0 \in \overline{C}\)。假设 \(R\) 不可约，并令 \(\beta\) 为 \(R\) 的最高根（关于 \(C'\)）；则 \(C\) 是 \(x^{*} \in C'\) 中满足 \(\langle x^{*}, \beta \rangle < 1\) 的元素所成的集合。

\begin{enumerate}
\def\labelenumi{\arabic{enumi})}
\setcounter{enumi}{31}
\item
  令 S 为 V 的对称代数。\(S^{W}\) 是由在 \(W = W(R)\) 下不变的元素组成的子代数，g 为 W 的阶。\(l = \dim V\)。存在齐次的、代数独立的元素 \(I_{1}, I_{2}, \ldots, I_{l}\)，它们生成 \(S^{W}\)。S 作为 \(S^{W}\)-模有一个由 g 个齐次元素组成的基。令 a 为 S 中由 \(S^{W}\) 的次数 \textgreater0 的齐次元素生成的理想：W 在 S/a 上的表示，是 W 在 S 上的表示通过取商所诱导的表示，同构于 W 的正则表示（在 R 上）。
\item
  令 \(I_1, I_2, \ldots, I_l\) 为生成 \(S^W\) 的齐次代数独立元素。其次数 \(k_1, k_2, \ldots, k_l\) 由 \(R\) 唯一确定（除排列外）。我们有 \(g = k_1k_2\ldots k_l\)。根的数目为 \(2\sum_{i=1}^{l}(k_i - 1)\)。
\item
  S 中的元素 A 若满足 \(w(A) = \det(w)\) 对所有 \(w \in W\) 成立，则称其在 W 下反不变。令 \(R = R_{1} \cup (-R_{1})\) 为 R 的一个划分，并置 \(\pi = \prod_{\alpha \in R_{1}} \alpha\)。S 中的元素 \(\pi\) 是反不变的：S 的反不变元素正是 \(\pi I\) 的形式，其中 \(I \in S^{W}\)。
\item
  令 E 为 R 的权群 P 的群代数 Z[P]。若 \(p \in P\)，我们用 \(e^{p}\) 表示 E 中相应的元素。我们有 \(e^{p}e^{p'} = e^{p+p'}\)，并且 \(e^{p}\) 构成 E 的一个基。W 在 E 上有一个作用，使得 \(w(e^{p}) = e^{w(p)}\) 在 \(w \in W\) 且 \(p \in P\) 时成立。若 E 中的元素 \(z \in E\) 满足 \(w(z) = \det(w).z\) 对所有 \(w \in W\) 成立，则称其为反不变元素。对所有 \(z \in E\)，置 \(J(z) = \sum_{w \in W} \det(w).w(z)\)。令 C 为一个室。\(J(e^{p})\)（其中 \(p \in P \cap C\)）构成 E 的反不变元素群的一个基。若 \(\rho\) 是正根之和的一半，则有：
\end{enumerate}

\[
J \left(e ^ {\rho}\right) = e ^ {\rho} \prod_ {\alpha > 0} \left(1 - e ^ {- \alpha}\right) = \prod_ {\alpha > 0} \left(e ^ {\alpha / 2} - e ^ {- \alpha / 2}\right),
\]

这些乘积是对根 \textgreater{} 0 的集合取的。

\begin{enumerate}
\def\labelenumi{\arabic{enumi})}
\setcounter{enumi}{35}
\tightlist
\item
按 35) 的记号，置
\end{enumerate}

\[
z _ {p} = \mathrm{J} (e ^ {p + \rho}) / \mathrm{J} (e ^ {\rho}) \quad \text { for } p \in \mathrm{P}.
\]

\(z_{p}\)（其中 \(p \in \mathrm{P} \cap \overline{\mathrm{C}}\)）构成由 W 不变的 E 中元素组成的群 \(\mathbf{E}^{\mathbf{W}}\) 的一个基。若 \(\omega_{1}, \ldots, \omega_{l}\) 是 R 的基本权，则元素 \(z_{\omega_{i}}\)（\(1 \leqslant i \leqslant l\)）代数独立并生成环 \(\mathbf{E}^{\mathbf{W}}\)。

\begin{enumerate}
\def\labelenumi{\arabic{enumi})}
\setcounter{enumi}{36}
\tightlist
\item
  令 C 为 R 的一个室，\((\alpha_{1},\ldots,\alpha_{l})\) 为相应的基。W 中的元素 \(c=s_{1}s_{2}\ldots s_{l}\) 称为 R 的 Coxeter 变换。c 的共轭类不依赖于 C 或 \(\alpha_{i}\) 的编号。c 的阶 h 称为 R 的 Coxeter 数。c 的特征值形如 \(\exp\frac{2i\pi m_{j}}{h}\)，其中整数 \(m_{1},m_{2},\ldots,m_{l}\)（称为 R 的指数）满足 \(1\leqslant m_{1}\leqslant m_{2}\leqslant\cdots\leqslant m_{l}\leqslant h-1\)。
\end{enumerate}

假设 R 不可约。则

\[
m _ {1} = 1, m _ {l} = h - 1.
\]

\[
m _ {j} + m _ {l + 1 - j} = h \quad (1 \leqslant j \leqslant l).
\]

\[
m _ {1} + m _ {2} + \dots + m _ {l} = \frac {1}{2} l h = \frac {1}{2} \operatorname{Card} (\mathrm{R}).
\]

每个与 h 互素的 \(m \in \{1, 2, \ldots, h - 1\}\) 恰好等于某一个 \(m_j\)。数 \(m_1 + 1, m_2 + 1, \ldots, m_l + 1\) 与 33) 中记为 \(k_1, k_2, \ldots, k_l\) 的整数重合（除排列外）。按 25) 的记号，\(n_1 + \cdots + n_l = h - 1\)。在 R 上，\(\{1, c, c^2, \ldots, c^{h-1}\}\) 有 l 个轨道，每个轨道含有 h 个元素。

若 \(h\) 为偶数，则 \(c^{h/2}\) 将 \(C\) 变为 \(-C\)。我们有 \(-1 \in W\)，当且仅当 \(W\) 的指数全为奇数；在这种情况下，\(h\) 为偶数且 \(c^{h/2} = -1\)。

\pandocbounded{\includegraphics[keepaspectratio]{images/part002_b04bf376.jpg}}
